\section{Transport of Q-zero faces through certified Pachner moves}
\label{sec:q-zero-pachner-transport}

A certified Pachner replacement preserves more than the underlying manifold
on a particular coordinate face of the normal matching cone.  If every
quadrilateral coordinate in the replaced region is zero, its triangle
coordinates are determined by five formal vertex weights.  Copying those
weights gives an exact correspondence that preserves canonical lifting,
formal Euler characteristic and root-ray extremality.  The argument allows
arbitrary quadrilateral coordinates outside the region, including
incompatible ones.  It therefore applies to the formal cones used by the
Euler search, rather than only to embedded normal surfaces.

The relative-boundary premises are those established in
Section~\ref{sec:strict-pachner32}.  The result applies to each elementary
move in the geometric descent of
Section~\ref{sec:geometric-descent-crushing}, provided its own regional
Q-zero condition holds.  It is a local transport theorem.  It does not
compare the full resolving parameters of arbitrary retriangulations.

\subsection{The checked formal region and its complete attachments}

Let $T_2$ and $T_3$ be validated finite connected orientable triangulations
with one torus boundary component, related by an accepted elementary
$2\leftrightarrow3$ certificate.  No one-vertex or knot-exterior hypothesis
is imposed.  Write their formal regions as
\[
\begin{aligned}
 R_2&:\quad [0,1,2,3],\ [0,1,2,4],\\
 R_3&:\quad [3,4,0,1],\ [3,4,1,2],\ [3,4,2,0].
\end{aligned}
\]
The five labels are distinct before attachment.  Equal formal internal
faces are paired by the maps preserving their vertex labels.  The internal
face of $R_2$ is $[0,1,2]$; those of $R_3$ are $[3,4,i]$ for $i=0,1,2$.
Both regions have the same six formal boundary faces
\[
 [3,0,1],\ [3,1,2],\ [3,2,0],\qquad
 [4,0,1],\ [4,1,2],\ [4,2,0].
\]

Let $O$ denote the corresponding outside tetrahedra.  Their local vertex
labels and mutual face maps are unchanged under the certified survivor
renumbering.  Every formal boundary face retains its complete attachment
record: it is unpaired, attached to a specified outside face by a specified
corner map, or attached to another formal boundary face by a specified
corner map.  Equality of face names alone would not suffice.  These full
records are exactly the data compared by the independent elementary
checkers described in Section~\ref{sec:strict-pachner32}.

The formal boundary may have selfattachments.  Different formal vertices
may represent the same global vertex, and different formal boundary edges
may represent the same global edge.  The proof does not require the image
of the closed bipyramid to be embedded in the quotient.  The outside set
$O$ may be empty.

For $i=2,3$, let $M_i$ be the original STANDARD matching matrix of $T_i$,
with four triangle and three Q coordinates per tetrahedron.  Define
\[
 \mathcal V_i^0=\{u:M_i u=0,\ q|_{R_i}=0\},\qquad
 \mathcal N_i^0=\mathcal V_i^0\cap\{u\ge0\}.
\]
Here $q|_{R_i}=0$ requires all three Q coordinates to vanish in every
tetrahedron of the indicated region.  The spaces $\mathcal V_i^0$ allow
signed real coordinates.  Neither definition imposes quadrilateral
compatibility in the outside tetrahedra.

\begin{proposition}[Exact transport on the regional Q-zero face]
\label{prop:q-zero-pachner-transport}
There is a linear isomorphism
$\Phi:\mathcal V_2^0\longrightarrow\mathcal V_3^0$ that copies every
outside STANDARD coordinate and replaces the regional triangles by the
same five formal vertex weights.  It has the following properties.
\begin{enumerate}
\item It identifies the nonnegative cones and the integral matching
      lattices.  The gcd of an integral vector is unchanged.
\item Global vertex classes correspond, their triangle-link vectors
      correspond, and $\Phi$ commutes with canonicalization at every
      global vertex.
\item The signed formal STANDARD Euler functional is preserved,
      including when outside Q coordinates are incompatible.
\item After identifying outside Q coordinates, the two regional Q-zero
      faces of the full Q cones are identical.  Nonzero root Q extreme
      rays in these faces correspond, with unchanged canonical formal
      Euler value, primitive Q scale and conflict sets under survivor
      relabelling.
\end{enumerate}
\end{proposition}

\begin{proof}
In a Q-zero tetrahedron, the arc on a face cutting off a local vertex
has weight equal to that vertex's triangle coordinate.  Internal face
matching in $R_2$ equates the two occurrences of each of $0,1,2$; vertices
$3,4$ occur once each.  Internal matching in $R_3$ equates the two
occurrences of each of $0,1,2$ and connects the three occurrences of each
of $3,4$.  Thus either regional matching space is exactly a five-weight
space with coordinates $a_0,a_1,a_2,a_3,a_4$.

In the displayed local corner orders, the map is
\[
\begin{array}{c|c}
 R_2 & R_3\\ \hline
 (a_0,a_1,a_2,a_3;0,0,0)&(a_3,a_4,a_0,a_1;0,0,0)\\
 (a_0,a_1,a_2,a_4;0,0,0)&(a_3,a_4,a_1,a_2;0,0,0)\\
                         &(a_3,a_4,a_2,a_0;0,0,0).
\end{array}
\]
The certified formal label maps convert this display to actual local
corner order.  Each $a_j$ can be read from any of its occurrences;
internal matching ensures they agree.  Copy the outside rows using the
survivor correspondence.  Reading the same five weights in the reverse
direction defines the inverse map.

On every formal boundary face, the arc at formal corner $j$ has weight
$a_j$ on both sides.  If the face is attached to an outside face $(u,g)$
and sends $j$ to local corner $v$, its matching equation on either side is
\[
 a_j=\tau_{u,v}+q_{u,\operatorname{quad}(g,v)},
\]
where $\tau_{u,v}$ denotes the outside triangle coordinate and
$\operatorname{quad}(g,v)$ is the type disjoint from edge $\{g,v\}$.
Both terms on the right and the corner map are unchanged.  Other Q
coordinates of that outside tetrahedron may be nonzero.  A boundary
selfattachment sending $j$ to $k$ instead gives $a_j=a_k$ on both sides.
Unpaired faces impose no matching equation.  Hence all remaining matching
equations are identical, establishing the linear isomorphism.

Every formal weight occurs on both sides, and every outside coordinate
is copied.  A vector is therefore nonnegative on one side exactly when
its image is nonnegative on the other.  The set of coordinate values is
unchanged, although their repetitions can change.  Integral coordinates
remain integral in both directions, and their gcd is unchanged.

To compare global vertices, start with the five formal labels and all
outside local corners.  The internal region maps identify only occurrences
of the same formal label.  All five labels occur on the formal boundary.
The remaining equivalence relation is generated by the identical boundary
attachments and outside face maps.  Thus global vertex classes correspond.
The new interior edge $[3,4]$ does not itself identify its endpoints as
vertices.  Such an identification, if present, already belongs to the
common attachment relation.

For a corresponding global vertex $v$, the set of triangle values at its
corners consists of its unchanged outside values together with
\[
                  \{a_j:j\text{ represents }v\}.
\]
Every such $a_j$ occurs on both sides.  Consequently the minimum $m_v$
is unchanged.  Write $L_v$ for the triangle-link vector equal to one at
every corner of $v$ and zero elsewhere.  Then $\Phi(L_v)=L'_v$, and
\[
 \Phi\left(u-\sum_v m_v L_v\right)
       =\Phi(u)-\sum_v m_v L'_v.
\]
This proves commutation with canonicalization.  It applies to signed
lifts as well as nonnegative ones, and includes interior vertices.
It requires the minimum at every global vertex, rather than a single
minimum over all corners.

The Euler calculation can be made in the signed matching space.  The
formal functional is
\[
 \chi(u)=\sum_{\text{global edges }e} w_e(u)
       -\sum_{\text{global faces }f} b_f(u)
       +\sum_{\text{tetrahedra }t} d_t(u),
\]
where $w_e$ is the edge-intersection weight, $b_f$ is the total face-arc
weight, and $d_t$ is the sum of the seven disc coordinates.  A paired face
is counted once, and an unpaired face is also counted once.  Edge weights
are well defined even for signed or incompatible coordinates: on a face,
the intersection weight at an edge is the sum of the two arc types at
its endpoints.  The original arc-matching equations therefore equate
edge weights under every face map.

Put $s=a_3+a_4$.  The regional disc counts before and after are
\[
 2(a_0+a_1+a_2)+s,\qquad 2(a_0+a_1+a_2)+3s,
\]
so the increase in disc count is $2s$.  The old internal face has arc
count $a_0+a_1+a_2$; the three new internal faces have total arc count
$a_0+a_1+a_2+3s$.  Thus the increase in internal arc count is $3s$.
The formal boundary traces, their quotient face pairings and the outside
face contributions are unchanged.

The nine formal boundary edges likewise have identical weights
$a_j+a_k$ and identical boundary attachment relations on both sides.
Their regional occurrences are identified by the internal maps to
precisely these same nine formal edges.  Applying the unchanged attachment
relations therefore gives corresponding old global edge classes with
unchanged weights.  The sole extra edge is $[3,4]$, of weight $s$.
Its three regional occurrences are paired only to one another by the
internal faces.  No formal boundary face contains this edge, so its
orbit cannot reach an outside or formal boundary edge.  This remains
true when $3$ and $4$ represent the same global vertex.  The change in
global edge-intersection weight is exactly $s$.  It follows that
\[
                   \chi(\Phi(u))-\chi(u)=s-3s+2s=0.
\]
This is an identity of linear functionals on signed matching vectors;
it uses no embedded-surface interpretation for the outside coordinates.

It remains to pass to Q coordinates.  Let $K_i$ be the full nonnegative
Q matching cone of $T_i$, and put
\[
                 P_i=\{q\in K_i:q|_{R_i}=0\}.
\]
Triangle-only matching solutions are exactly the span of the global
vertex-link vectors.  Every signed matching lift of a nonnegative Q
vector becomes nonnegative after adding sufficient multiples of these
links, or after canonicalization.  Hence $K_i$ is also the projection
of the nonnegative STANDARD matching cone.  The maps $\Phi$ and
$\Phi^{-1}$ supply lifts in both directions and copy every outside Q
coordinate.  After deleting regional zeros, $P_2$ and $P_3$ are therefore
the same subset of $\mathbb R^{3|O|}$.

Both $P_i$ are coordinate faces of their full cones.  If a nonnegative
vector with the required zeros decomposes in $K_i$, both summands have
the same zeros.  Thus a nonzero vector is extreme in $P_i$ exactly when
it is extreme in $K_i$, as in Lemma~\ref{lem:transversal-face-rays}.
The face identity proves the asserted root-ray correspondence.
Commutation with canonicalization and preservation of $\chi$ give equal
canonical formal Euler values.  Outside Q entries are copied and all
regional entries are zero, so Q gcd and conflicts are unchanged.  This
also proves preservation of compatibility or incompatibility.
\end{proof}

Formal vertex aliasing must not be replaced by equality of triangle
weights in this proof.  If two formal vertices represent the same global
vertex, an outside corner path can change triangle potentials through
nonzero Q terms.  Only actual matching equations impose equalities among
the $a_j$.  Canonicalization uses their unchanged set of values within
each global vertex class, without requiring those values to be equal.

\subsection{Original-matrix rank and primitive canonical lifts}

The cone statement has a rank formulation using the original STANDARD
rows, without trusting a particular triangle elimination.  This is the
same source-equation perspective used for the independent kernel checks
of Section~\ref{sec:reusable-source-kernels}.

\begin{corollary}[Supported Q nullity and the four-column rank change]
\label{cor:q-zero-pachner-rank}
Fix a set $S$ of outside Q coordinates.  On each side retain all triangle
columns and precisely the Q columns in $S$ in the original matching
matrix; call the resulting matrix $A_i(S)$.  Then
\[
 \dim\ker A_2(S)=\dim\ker A_3(S),\qquad
 \operatorname{rank}A_3(S)=\operatorname{rank}A_2(S)+4.
\]
The corresponding all-triangle ranks also differ by four, and the
supported projected Q nullities are equal.  In particular a nonzero Q
vector positive on exactly $S$ passes the original-matrix root-ray
nullity-one test on one side exactly when it does on the other.
\end{corollary}
\begin{proof}
Restricting $\Phi$ to zero outside Q coordinates off $S$ identifies the
two signed kernels.  A $2\to3$ move adds four triangle columns and no
retained Q columns.  Equal nullities therefore give the rank difference
four.  The same argument with $S=\varnothing$ gives the all-triangle
claim.  Projection to supported Q coordinates has as kernel the span
of the global vertex links, which correspond under $\Phi$.  Quotienting
the two isomorphic signed kernels by these subspaces proves equality
of projected Q nullities.  At a Q vector positive on its entire support,
nullity one is precisely the extreme-ray condition: a higher-dimensional
kernel permits sufficiently small nonproportional perturbations in both
directions while preserving nonnegativity.
\end{proof}

The rank formula retains \emph{all triangle columns}.  It does not state
that the number or rank of positive STANDARD support columns of a
particular canonical vector changes by four.  Those support columns can
change according to which formal weights vanish.

For an integral Q vector, the canonical STANDARD lift is integral.
Indeed each triangle matching equation is a triangle difference plus an
integral combination of Q coordinates.  Propagating potentials from a
zero root along a corner spanning forest gives integral triangle values;
Q matching supplies consistency around the remaining cycles.  Subtracting
an integral minimum at each global vertex preserves integrality.  Since
the Q coordinates themselves are copied into the lift, primitive integral
Q gives a primitive STANDARD lift.  This is the integral lifting argument
of Section~\ref{sec:direct-crush-positive}; the transport additionally
preserves the full STANDARD gcd by copying the same set of values.

\subsection{Sequences, protected conflicts and the boundary of the theorem}

For a certified sequence, apply
Proposition~\ref{prop:q-zero-pachner-transport} at each intermediate
source.  In particular a geometric $UDD$ descent requires regional
Q zeros for its expansion and for each of its two contractions, using
the vector already transported to that step.  The two contraction
regions can differ from the expansion region.  Checking only the
initial or final source is insufficient.

The condition concerns complete Q support.  Having at most one positive
Q type in a moved tetrahedron, having no conflict there, or having
disjoint conflict sets for several rays does not establish regional
Q zeros.  If there are no outside tetrahedra, the protected Q face is
the zero cone; the signed matching and vertex-link statements remain
valid, but there is no protected nonzero Q ray.

A family of positive incompatible root rays satisfying the condition
at every step transports with the same conflict hypergraph after
relabeling its surviving tetrahedra.  Any proved transversal lower
bound for that family therefore persists.  This does not identify
the unrestricted cones or constrain their other rays.  It gives no
equality, monotonicity, or additive bound for the full-source resolving
numbers $h(T)$ of Section~\ref{sec:euler-transversal}.  A negative proof
on a protected coordinate face does not certify the unrestricted
successor cone, and old LP dual multipliers cannot simply be copied
between different source matrices and gauges.  Transport through a
normal crush requires a separate correspondence theorem.

To obtain a growing resolving-parameter obstruction from this local
result, one would still need an actual family of the intended
knot-exterior sources, sufficiently many positive root rays with
provably large conflict transversal number, protection through every
chosen move, and stationarity under the complete preprocessing policy.
Unbounded growth would exclude a constant bound; excluding a
polylogarithmic bound requires correspondingly stronger growth relative
to source size.  None of those global construction obligations follows
from the local theorem.

\subsection{Implementation premises and prospective reuse}

The proof's local maps agree with the maintained independent
\path{verify_pachner_23} and \path{verify_pachner_32} checkers: full formal
boundary signatures include region-to-region corner maps, and both
sources undergo finite manifold validation.  Its canonicalization is
the per-global-vertex minimum used by the source kernel.  Its signed
Euler calculation is the linear extension of the original edge, face
and disc census used by the independent Euler model.  These code
premises were reviewed separately from the theorem's algebra; they
supply no additional empirical or performance claim.

A future transport adapter could authenticate the old witness and
elementary move, compare the five formal weights, copy outside rows,
and fill the new region.  For $b$-bit integral coordinates, the copying
and comparisons require $O(t(b+\log t))$ bit work, including label
bookkeeping but apart from source and packet
validation.  A conservative first adapter could still replay matching,
canonical Euler and root rank independently on the successor.  Avoiding
that repeated rank computation would require a dedicated checked
transport contract.  Such an adapter must bind both labelled sources,
respect the exact survivor and corner maps, and shield caller callback
exceptions across the legacy checker calls.

No native transport adapter or persistent witness cache is introduced
here.  The current driver completes its polynomial preprocessing before
Euler discovery, and a normal crush separates later searches.  It
therefore usually has no witness from before a Pachner move to reuse.
Persistent source-bound witnesses or a changed interleaving policy
would be needed to exploit the theorem in those runs.  Preserving
incompatible rays can also preserve search obstructions, so reuse or
move ordering has no automatic performance benefit.  This section
proves the local correspondence and leaves those algorithmic uses and
the general quasi-polynomial recognition bound open.

There is also no immediate reduction of the current branch LP dimension.
Starting with the full mask, both search policies retain at least one
Q type in every tetrahedron: a split retains a singleton and includes
the zero-Q possibility within that singleton cone.  A zero regional
support in one returned point does not impose zeros on its whole node.
Even on a supplied mask that does impose regional zeros, the theorem
identifies the same outside Q variables and the same linear matching
space.  Restricted triangle assembly already contracts zero-labelled
triangle equations, and original-row selection removes redundant Q
equations.  The rank-reduced LP therefore has the same intrinsic equality
rank before and after transport.  Repeated matrix preparation across
different objectives of that LP is a separate algebraic cost, addressed
in Section~\ref{sec:prepared-euler-matrix}.
